\section{Métricas de Centralidad}

\subsection{Degree Centrality}

\textbf{Concepto.} Un vértice es importante si tiene muchas conexiones directas.
Es la métrica más simple y sirve como referencia para comparar con las demás.

\textbf{Fórmula} (grafo no dirigido, normalizada):
\[
C_D(v) = \frac{\deg(v)}{n - 1}
\]
donde $\deg(v)$ es el número de vecinos directos de $v$ y $n = |V|$.

En grafos \textbf{dirigidos} se distingue in-degree (aristas entrantes, popularidad)
y out-degree (aristas salientes, actividad).

\textbf{Complejidad:} $O(|V| + |E|)$

\textbf{Aplicación.} En la red Yeast PPI, una proteína con alto degree es un
\textit{hub} biológico cuya eliminación puede resultar letal para el organismo.
En Trade Network, el in-degree refleja diversificación de importaciones y el
out-degree el alcance exportador de un país.

% ────────────────────────────────────────────────
\subsection{Betweenness Centrality}

\textbf{Concepto.} Un vértice es importante si actúa como puente en los caminos
más cortos entre otros pares. Captura el control sobre el flujo en la red.

\textbf{Fórmula:}
\[
C_B(v) = \sum_{\substack{s \neq v \neq t \\ s \neq t}}
         \frac{\sigma_{st}(v)}{\sigma_{st}}
\]
donde $\sigma_{st}$ es el número total de caminos más cortos entre $s$ y $t$,
y $\sigma_{st}(v)$ los que pasan por $v$.

\textbf{Normalización:}
\[
C_B^{\text{norm}}(v) =
\begin{cases}
\dfrac{C_B(v)}{(n-1)(n-2)} & \text{dirigido} \\[8pt]
\dfrac{2\,C_B(v)}{(n-1)(n-2)} & \text{no dirigido}
\end{cases}
\]

\textbf{Algoritmo:} Brandes~\cite{brandes2001}, que calcula la betweenness exacta
para todos los vértices con un BFS por vértice fuente.

\textbf{Complejidad:} $O(|V| \cdot |E|)$

\textbf{Aplicación.} En Yeast PPI, alta betweenness indica proteínas que conectan
módulos funcionales distintos; su eliminación puede desconectar regiones de la red.
En Trade Network, identifica países que actúan como intermediarios críticos del
comercio global.

% ────────────────────────────────────────────────
\subsection{Closeness Centrality}

\textbf{Concepto.} Un vértice es importante si puede alcanzar rápidamente a todos
los demás. Se basa en la inversa de la distancia promedio.

\textbf{Fórmula} (Wasserman \& Faust~\cite{wasserman1994}, para grafos desconectados):
\[
C_C(v) = \frac{n-1}{N-1} \cdot \frac{n-1}{\displaystyle\sum_{u \neq v} d(v,u)}
\]
donde $n$ es el número de vértices alcanzables desde $v$, $N = |V|$ el total,
y $d(v,u)$ la distancia mínima entre $v$ y $u$.

El factor $\frac{n-1}{N-1}$ penaliza vértices que no alcanzan al resto, permitiendo
comparar en grafos desconectados.

\textbf{Complejidad:} $O(|V| \cdot (|V|+|E|))$ via BFS desde cada vértice.

\textbf{Aplicación.} En Yeast PPI, alta closeness indica proteínas con rol
regulatorio global, capaces de transmitir señales al resto de la red en pocas
interacciones. En Trade Network, países con alta closeness están integrados
eficientemente con la economía mundial.

% ────────────────────────────────────────────────
\subsection{PageRank}

\textbf{Concepto.} Un vértice es importante si está conectado a otros vértices
importantes. Un enlace desde un nodo de alto PageRank vale más~\cite{pagerank1999}.

\textbf{Fórmula iterativa:}
\[
PR(v) = \frac{1-d}{n} + d \sum_{u \in N^{-}(v)} \frac{PR(u)}{|N^{+}(u)|}
\]
donde $d = 0.85$ es el factor de amortiguación (\textit{damping factor}),
$N^{-}(v)$ son los vecinos que apuntan a $v$, y $N^{+}(u)$ los destinos de $u$.

El término $\frac{1-d}{n}$ modela la probabilidad de salto aleatorio a cualquier nodo.
El algoritmo itera hasta convergencia ($\|PR^{(k+1)} - PR^{(k)}\|_1 < 10^{-6}$).

\textbf{Complejidad:} $O(k \cdot (|V|+|E|))$, con $k < 100$ en la práctica.

\textbf{Aplicación.} En Trade Network, el PageRank refleja la importancia de un
país en el comercio global ponderada por la importancia de sus socios, capturando
influencia indirecta que el degree no detecta. En Yeast PPI, identifica proteínas
centrales cuya importancia está reforzada por sus vecinas.

% ────────────────────────────────────────────────
\subsection{Average Shortest Path}

\textbf{Concepto.} Promedio de las distancias mínimas desde un vértice hacia todos
los demás alcanzables. Un valor alto indica aislamiento estructural.

\textbf{Fórmula por vértice:}
\[
ASP(v) = \frac{1}{|R(v)|} \sum_{u \in R(v)} d(v, u)
\]
donde $R(v)$ son los vértices alcanzables desde $v$ (excluyendo $v$).

\textbf{Métrica global:}
\[
\overline{ASP} = \frac{1}{|V|} \sum_{v \in V} ASP(v)
\]
Esta métrica global cuantifica la compacidad de la red entera. El experimento de
los \textit{seis grados de separación} corresponde a medir $\overline{ASP}$ en la
red social humana.

\textbf{Complejidad:} $O(|V| \cdot (|V|+|E|))$

\textbf{Aplicación.} En Yeast PPI, un $\overline{ASP}$ bajo confirma que la red
es compacta (propiedad de \textit{small-world}), lo que facilita la propagación
de señales biológicas. En Trade Network, países con bajo $ASP$ individual son los
más integrados al sistema de comercio global.

% ────────────────────────────────────────────────
\subsection{Eigenvector Centrality}

\textbf{Concepto.} Extiende el degree ponderando cada vecino por su propia
centralidad. Es la base matemática del PageRank~\cite{bonacich1987}.

\textbf{Definición.} Sea $\mathbf{A}$ la matriz de adyacencia. El vector
$\mathbf{x}$ de Eigenvector Centrality satisface:
\[
\mathbf{A}\,\mathbf{x} = \lambda\,\mathbf{x}
\qquad \Longleftrightarrow \qquad
C_E(v) = \frac{1}{\lambda} \sum_{u \in N(v)} C_E(u)
\]
donde $\lambda$ es el autovalor dominante de $\mathbf{A}$ (Perron-Frobenius).

\textbf{Algoritmo:} \textit{Power Iteration}: inicializar $\mathbf{x}^{(0)}$
uniformemente, iterar $\mathbf{x}^{(k+1)} \leftarrow \mathbf{A}\,\mathbf{x}^{(k)}$
normalizando por el máximo hasta convergencia.

\textbf{Complejidad:} $O(k \cdot (|V|+|E|))$

\textbf{Diferencia con PageRank.} PageRank agrega un factor de amortiguación $d$ y
maneja nodos sin salidas. Eigenvector Centrality es la versión pura; ambos coinciden
cuando $d \to 1$ y no hay nodos dangling.

\textbf{Aplicación.} En Yeast PPI, identifica el núcleo funcional: proteínas que
interactúan con otras que también son muy conectadas. En Trade Network, captura países
cuya influencia comercial se amplifica por sus socios.

% ────────────────────────────────────────────────
\subsection{Clustering Coefficient}

\textbf{Concepto.} Mide qué fracción de los vecinos de un vértice están conectados
entre sí. Cuantifica la tendencia a formar grupos locales (\textit{clustering})~\cite{watts1998}.

\textbf{Fórmula local:}
\[
C_{Cl}(v) = \frac{|\{(u,w) \in E \;:\; u,w \in N(v)\}|}{\deg(v)\,(\deg(v)-1)}
\]
El numerador cuenta los pares de vecinos de $v$ que están conectados (triángulos
que pasan por $v$); el denominador es el máximo posible. El resultado está en $[0,1]$.

\textbf{Coeficiente global:}
\[
\overline{C_{Cl}} = \frac{1}{|V|} \sum_{v \in V} C_{Cl}(v)
\]

\textbf{Complejidad:} $O(|V| \cdot \overline{\deg}^{\,2})$, donde $\overline{\deg}$ es el grado promedio.

\textbf{Aplicación.} En Yeast PPI, alto clustering indica la presencia de complejos
proteicos: grupos donde todas las proteínas interactúan entre sí (ribosomas, complejos
de transcripción). En Trade Network, alto clustering de un país indica que sus socios
comerciales también comercian entre sí, formando bloques económicos regionales.
